\vspace{-0.1in}
\subsection{Implementation Setup}
\textbf{Baselines. }We evaluate the zero-shot performance of 4 recently released LMMs: (1) MiniGPT4 (7B) \cite{zhu2023minigpt}; (2) LLaVA (7B) \cite{liu2023visual}; (3) InstructBLIP (7B) \cite{dai2023instructblip}; (4) Multimodal-GPT (7B) (MMGPT) \cite{gong2023multimodal}. All the above models have been tuned on their own collected visual instruction data. 


\textbf{Training Details. }We initialize the model from the first pretraining stage checkpoint of MiniGPT4. Then we instruct-tune the model on our dataset with the linear projection layer as the only learnable module. As for the hyper-parameters, we  employ a batch size of 1, a linear warmup of the learning rate 1e-6 during the initial 200 steps, and a weight decay of weight decay 0.05. We trained our models on NVIDIA Quadro RTX 8000 %8 Nvidia A100 (48G) 
GPUs with the maximum epoch as 20. 
% In addition, 
We randomly select 1000 instances from our dataset for evaluation, and the others are used for finetuning.


% \kevin{The below question list can be removed. Repeated in Sec6.2 already}

% \textbf{Research Questions. }We conduct empirical experiments to investigate the following questions: 
% \begin{itemize}[nosep,leftmargin=*]
%     \item \textit{Question1: "Do LMMs perform better on Positive or Negative instructions?"}. 
%     \item \textit{Question2: "Do LMMs perform better on different formats and lengths of instructions?"}. 
%     \item \textit{Question3: "How do LMMs perform on different semantic levels of hallucination?"}
%     \item \textit{Question4: "How do LMMs perform on different datasets?"}
%     \item \textit{Question5: "How do LMMs perform on different ratios of composition in training data?"} 
% \end{itemize}


\subsection{Main Results}
% \textbf{Results on \textit{LRV-Instruction}. }
\textbf{\textit{How do LMMs perform on \textit{LRV-Instruction}? }}We show the evaluation results on our dataset in Tab.~\ref{tab:comparison_evaluation}. Among the baselines, InstructBLIP achieves better results than other LMM baselines because its visual instructions are collected from a wide variety of publicly available datasets. LLaVA \cite{liu2023visual} utilizes the GPT-assisted approach to generate visual instructions, but its performance is much worse. This is probably because its synthetic answers from GPT4 are generally longer and may involve irrelevant information. As a comparison, our model outperforms the existing LMM baselines by a large margin, which benefits from the rich composition of our dataset and better prompt design. %\kevin{I would suggest to have a subsection "Main Result" at the beginning of Sec 6.2, and put this paragraph in the subsection. After that, add the title "How do LMMs perform on different datasets?" to the next paragraph. }
%\kevin{Another way: Change "How do LMMs perform on different datasets?" to "Effectiveness of LRV-Instruct..."} \kevin{Another way: at the beginning, mention we test on three datasets, including AAA,BBB,CCC. Then discuss each result.}


% \begin{table}[t]
% \begin{subtable}[t]{0.5\textwidth}
% \centering
% \small
% \begin{tabular}{lcc}
% \toprule[1.5pt]
% Model & Training Size & Accuracy\\
% \midrule
% InstructBLIP & >1M & \textbf{0.62}\\
% LLaVA & 150k & 0.48\\
% MiniGPT4 & 3k & 0.44\\
% Ours & 120k& \textbf{0.61}\\
% \bottomrule[1.5pt]
% \multicolumn{1}{c}{} \\ % space holder
% \end{tabular}
%  \caption{Evaluation results on GQA dataset \cite{hudson2019gqa}.}
% \label{tab:perform11}
% \end{subtable}
% \hspace{\fill}
% \begin{subtable}[t]{0.5\textwidth}
% \centering
% \small
% \begin{tabular}{lcc}
% \toprule[1.5pt]
% Ratio & \textit{$Acc_p$} & \textit{$Acc_n$}\\
% \midrule
% \textit{All Pos} & 0.97 & 0.05\\
% \textit{Pos:Neg=2:1} & 0.95 & 0.50\\
% \textit{Pos:Neg=1:1} & 0.92 & 0.85\\
% \textit{Pos:Neg=1:2} & 0.87 & 0.86\\
% \textit{All Neg} & 0.10 & 0.98\\
% \bottomrule[1.5pt]
% \end{tabular}
% \caption{
% Results of different composition ratios in training.
% }
% \label{tab:perform112}
% \end{subtable}
% \vspace{-0.05in}
% \caption{
% Evaluation of current LMMs on the VQA dataset and ablation study of our model.
% }
% \label{tab:perform12}
% \vspace{-0.2in}
% \end{table}

%%%%%%%

\begin{table}[t]
\centering
\def\arraystretch{1.2}
\begin{minipage}{0.48\linewidth}\centering
\centering
\small
\begin{tabular}{lcc}
\toprule[1.5pt]
Model & Training Size & Accuracy\\
\midrule
InstructBLIP & >1M & \textbf{0.62}\\
LLaVA & 150k & 0.48\\
MiniGPT4 & 3k & 0.44\\
Ours & 120k& \textbf{0.61}\\
\bottomrule[1.5pt]
\multicolumn{1}{c}{} \\ % space holder
\end{tabular}
\vspace{-0.05in}
\caption{Results on GQA~\cite{hudson2019gqa}. We evaluate LMM's performance in general VQA scenario.}
\label{tab:perform11}
\end{minipage}
%
\begin{minipage}{0.48\linewidth}\centering
\begin{tabular}{lcc}
\toprule[1.5pt]
Ratio & \textit{$Acc_p$} & \textit{$Acc_n$}\\
\midrule
\textit{All Pos} & 0.97 & 0.05\\
\textit{Pos:Neg=2:1} & 0.95 & 0.50\\
\textit{Pos:Neg=1:1} & 0.92 & 0.85\\
\textit{Pos:Neg=1:2} & 0.87 & 0.86\\
\textit{All Neg} & 0.10 & 0.98\\
\bottomrule[1.5pt]
\end{tabular}
\vspace{0.1in}
\caption{Results of different composition ratios in visual instruction tuning. }
\label{tab:perform112}
\end{minipage}
\end{table}



% \begin{table*}[t!]\centering
% %%%%%%%%%% subfloat a
% \subfloat[Evaluation results on GQA dataset \cite{hudson2019gqa}\label{tab:perform11}]
% \begin{tabular}{lcc}
% 	\toprule[1.5pt]
% 	Model & Training Size & Accuracy\\
% 	\midrule
% 	InstructBLIP & >1M & \textbf{0.62}\\
% 	LLaVA & 150k & 0.48\\
% 	MiniGPT4 & 3k & 0.44\\
% 	Ours & 120k& \textbf{0.61}\\
% 	\bottomrule[1.5pt]
% \multicolumn{1}{c}{} \\ % space holder
% \end{tabular}
% \hfill
% %%%%%%%%%% subfloat b
% \subfloat[Results of different composition ratios in training \label{tab:perform112}]
% \begin{tabular}{lcc}
% 	\toprule[1.5pt]
% 	Ratio & \textit{$Acc_p$} & \textit{$Acc_n$}\\
% 	\midrule
% 	\textit{All Pos} & 0.97 & 0.05\\
% 	\textit{Pos:Neg=2:1} & 0.95 & 0.50\\
% 	\textit{Pos:Neg=1:1} & 0.92 & 0.85\\
% 	\textit{Pos:Neg=1:2} & 0.87 & 0.86\\
% 	\textit{All Neg} & 0.10 & 0.98\\
% 	\bottomrule[1.5pt]
% \end{tabular}
% \caption{
% Evaluation of current LMMs on the VQA dataset and ablation study of our model
% }
% \vspace{-4mm}
% \end{table*}



\begin{table}[t]
\setlength\tabcolsep{3pt}
\centering
\small
\begin{tabular}{lcccccc}
\toprule[1.5pt]
Categories& Metric&Ours & MiniGPT4 & LLaVA & InstructBLIP & MMGPT\\
\midrule
Neg1 & \textsc{Accuracy(GPT4)} & \textbf{8.90} & 3.72 & 2.09 & 5.50 & 1.13\\
Neg1 & \textsc{Relevancy(GPT4)} & \textbf{8.96} & 5.94 & 4.83 & 7.22 & 2.24\\
\midrule
Neg2 & \textsc{Accuracy(GPT4)} & \textbf{6.50} & 2.57 & 1.42 & 2.18 & 0.96\\
Neg2 & \textsc{Relevancy(GPT4)} & \textbf{8.46} & 2.53 & 1.82 & 2.73 & 1.19\\
\bottomrule[1.5pt]
\end{tabular}
\vspace{0.05in}
\caption{
Evaluation results on \textit{Neg1} instructions and \textit{Neg2} instructions by \textit{GAVIE}. The metric scores of \textsc{Accuracy(GPT4)} and \textsc{Relevancy(GPT4)} are in a scale of 0 to 10.}
\label{tab:neg1neg2}
\vspace{-0.1in}
\end{table}


\textbf{\textit{How do LMMs perform on different datasets? }}We compare our model against the baseline models on the hallucination evaluation dataset recently introduced by POPE~\cite{li2023evaluating}, which randomly selects 500 images from COCO \cite{lin2014microsoft}. Since the authors of POPE~\cite{li2023evaluating} have not released their dataset at the time of our paper submission, we employ our GPT4-assisted approach to generate open-ended negative instructions. In addition, we instruct the model to produce “Yes” or “No” and use classification accuracy and F1 for evaluation according to \cite{li2023evaluating}. As shown in Tab. \ref{tab:coco}, we observe that MiniGPT4, LLaVA and InstructBLIP may not work well with open-ended negative instructions. In contrast, the highest scores of our model demonstrates that our dataset is more suitable for robust visual instruction tuning. Although InstructBLIP tends to hallucinate less, its training set is much larger than ours (\textit{120k vs 1M}). Moreover, our model is built on MiniGPT4, which only instruct-tunes the linear projector layer. To further explore the LMMs' performance in the common scenario of visual question-answering (VQA), we follow \cite{li2023evaluating} to evaluate LMMs on GQA dataset~\cite{hudson2019gqa} in the zero-shot setting. It is worth noting that GQA consists of open-ended VQA questions. As shown in Tab. \ref{tab:perform11}, the results suggest that our method achieves on-par performance with InstructBLIP in generic VQA setting, and our method requires much less training data. %By finetuning MiniGPT4 on our dataset, it achieves on-par performance with InstructBLIP but employs much less training data.


%we follow their proposed method to construct the \textsc{Random Negative} instructions. To construct suchinstructions, we randomly sample the objects from 80 object classes in COCO that do not exist in the image. In addition, we instruct the model to produce “Yes” or “No” and use classification accuracy and F1 for evaluation. We adopt the same instruction templates from \cite{li2023evaluating}: \textit{(1) Is there an <object> in the image? (2) Does the image contain a <object>?} 

%As shown in Tab. \ref{tab:coco}, our model outperforms MiniGPT4 and LLaVA in the \textsc{Random Negative} instructions which demonstrates the effectiveness of our dataset in mitigating the hallucination issue. Although InstructBLIP tends to hallucinate less, its training set is much larger than ours (\textit{120k vs 1M}). Moreover, our model is built on MiniGPT4, which only instruct-tunes the linear projector layer. 

%In order to evaluate LMMs on non-template instructions in this dataset, we employ our GPT4-assisted approach to generate \textsc{Hard Negative} instructions. 

%From Tab.~\ref{tab:coco}, the overall performance of baselines is worse compared to the \textsc{Random Negative} setting, showing that baselines may not work well with open-ended negative instructions. In contrast, our model can achieve similar performance for both negative instruction categories, which demonstrates that our dataset is more suitable for robust visual instruction tuning. 





\subsection{Detailed Analysis}
\textbf{\textit{Do LMMs perform better on Positive or Negative Instructions? }}Our evaluation set consists of positive and negative instances. We divide it into two sets and analyze the model performance on each of them. As shown in Fig.~\ref{fig:posneg}, baseline models, including MiniGPT4, LLaVa, and InstructBLIP,  perform better on positive instances than negative ones, as the training data adopted by these models do not contain negative instructions. MMGPT has low performance on both sets due to many repetitive phrases in the response. In addition, we found that the degradation of LLaVA is the most severe. We hypothesize that the synthetic answers for instruction tuning in LLaVA are generally longer and involve more unrelated information. In contrast, our model performs the best in both sets. InstructBLIP performs with higher scores than other LMMs because of the effectiveness of its instruction-aware visual encoder to extract image information.



\textbf{\textit{How do LMMs perform at the different semantic levels of hallucination? }} As shown in Tab \ref{tab:neg1neg2}, all baselines perform better on \textit{Neg1} than \textit{Neg2}. From the visual perspective, existent object manipulations with wrong attributes in \textit{Neg2} are more challenging than adding nonexistent objects from images to instructions in \textit{Neg1}. For example, in Tab. \ref{fig:neg_example}, it %is not hard 
may be straightforward to find that the "hot air balloon" does not appear in the image. However, "woman" does exist in the second example of Fig.~\ref{fig:posneg} while she is not in the blue pants and pink shirts, which requires a fine-grained understanding of the visual content. Therefore, a more powerful vision encoder is needed for future LMMs. In contrast, our model improves a lot at both semantic levels, which benefits from our diverse instruction tuning data. In the Appendix, we present more visualizations of model predictions for \textit{Neg1} and  \textit{Neg2} examples. % Fig.~\ref{fig:demo_neg1} and \ref{fig:demo_neg2} in the appendix show more visualizations of model predictions for \textit{Neg1} and  \textit{Neg2} examples. %\linjie{Will Appendix be in the same file as the main text? If not, need to mention Figure \ref{fig:demo_neg1} and \ref{fig:demo_neg2} are in the appendix.}
% of different hallucinations.

% \kevin{we have discussed length and format in the previous paragraph. Here we should directly discuss observations at different semantic levels}




\textbf{\textit{How do LMMs perform at the different composition ratios in training data?}} Our model benefits from the datasets with both positive and negative instructions. In Tab. \ref{tab:perform112}, we investigate how it addresses hallucination issues with different ratios of positive and negative samples in the training set. Inspired by \cite{li2023evaluating}, we instruct the model to produce “Yes” or “No” and use classification accuracy for evaluation. \textit{$Acc_p$} is the accuracy on the positive instruction set, while \textit{$Acc_n$} is the accuracy on the negative instruction set. From Tab. \ref{tab:perform112}, we found that \textit{$Acc_n$} increases with more negative samples, which verifies our hypothesis that the hallucination problem of current LMMs \cite{zhu2023minigpt, liu2023visual, gong2023multimodal, dai2023instructblip} are due to lack of 
negative instructions. Besides, with a balanced ratio (\textit{pos:neg=1:1}), the model performs the best in both positive and negative sets. We also found that the setting \textit{pos:neg=1:2} can achieve a similar level of results as \textit{pos:neg=1:1}. This is probably because the backbone of our model (MiniGPT4) was trained on a large amount of image-caption pair datasets, which are positive samples.


\textbf{\textit{Do LMMs perform better on different formats and lengths of instructions? }}From Tab \ref{tab:Interrogative}, LMMs perform with higher scores on interrogative instructions than declarative, but the difference is relatively small. Even though recent visual instruction tuning datasets lack diverse declarative instructions, the LMMs built on LLM are powerful enough to understand and follow the declarative instructions. From Fig. \ref{fig:length}, current LMMs achieve better results in short instructions than long ones since longer instructions contain more information, making it more challenging to understand. 


\begin{figure}[t]
     \centering
     \begin{subfigure}[b]{0.49\textwidth}
         \centering
         \includegraphics[width=\textwidth]{sub1.pdf}
         \vspace{-0.2in}
         \caption{Accuracy Performance.}
         \label{fig:y equals x}
     \end{subfigure}
     \hfill
     \begin{subfigure}[b]{0.49\textwidth}
         \centering
         \includegraphics[width=\textwidth]{sub2.pdf}
         \vspace{-0.2in}
         \caption{Relevancy Performance.}
         \label{fig:three sin x}
     \end{subfigure}
        \caption{Evaluation results on positive and negative instructions by \textit{GAVIE}. }
        \label{fig:posneg}
\vspace{-0.1in}
\end{figure}

\begin{figure}[t]
     \centering
     \begin{subfigure}[b]{0.49\textwidth}
         \centering
         \includegraphics[width=\textwidth]{length1.pdf}
         \vspace{-0.2in}
         \caption{Accuracy Performance.}
         \label{length1}
     \end{subfigure}
     \hfill
     \begin{subfigure}[b]{0.49\textwidth}
         \centering
         \includegraphics[width=\textwidth]{length2.pdf}
         \vspace{-0.2in}
         \caption{Relevancy Performance.}
         \label{length2}
     \end{subfigure}
        \caption{Evaluation results on different instruction lengths by \textit{GAVIE}. }
        \label{fig:length}
\vspace{-0.2in}
\end{figure}


\begin{table}[h]
\centering
\small
\begin{tabular}{lcccccc}
\toprule[1.5pt]
Categories & Metric &Ours & MiniGPT4 & LLaVA & InstructBLIP & MMGPT\\
\midrule
Interrogative&\textsc{Accuracy(GPT4)} & \textbf{6.61} &  4.14 & 4.60 & 5.95 & 1.01\\
Interrogative&\textsc{Relevancy(GPT4)} & \textbf{8.46} &  6.20 & 5.88 & 7.67 & 2.00\\
\midrule
Declarative&\textsc{Accuracy(GPT4)} & \textbf{6.50} & 3.98 & 3.82 & 5.47 & 0.90\\
Declarative&\textsc{Relevancy(GPT4)} & \textbf{8.21} & 5.39 & 5.84 & 6.64 & 1.62\\
\bottomrule[1.5pt]
\end{tabular}
\vspace{0.05in}
\caption{
Evaluation results on Interrogative Instructions and Declarative Instructions by \textit{GAVIE}. The metric scores of \textsc{Accuracy(GPT4)} and \textsc{Relevancy(GPT4)} are in a scale of 0 to 10.
}
\label{tab:Interrogative}
\vspace{-0.2in}
\end{table}